\clearpage\pagestyle{otMain}\setlength{\otBlocEspacement}{2.5mm}\fontsize{11}{15}\selectfont\corrigesChapitre{11}{Proportionnalité et pourcentages}
\begin{blocLivre}{cor-11-1}\begin{corrige}{1}{Un coefficient}\demarche{Revenir à une unité.}
$7{,}50:3=\res{2{,}50}$ euros par kg. Le prix s’obtient en multipliant la masse par 2,5.
\end{corrige}\end{blocLivre}
\begin{blocLivre}{cor-11-2}\begin{corrige}{2}{Une recette doublée}\demarche{Multiplier les deux grandeurs par le même nombre.}
On double le nombre de personnes, donc la quantité : $180\times2=\res{360\ \mathrm g}$.
\end{corrige}\end{blocLivre}
\begin{blocLivre}{cor-11-3}\begin{corrige}{3}{Un pourcentage simple}\demarche{Utiliser les parts repères.}
On obtient $\res{24}$, $\res{120}$ et $\res{2{,}4}$. Ce sont respectivement un dixième, une moitié et un centième.
\end{corrige}\end{blocLivre}
\begin{blocLivre}{cor-11-4}\begin{corrige}{4}{Proportion ou effectif}\demarche{Rapporter la part au total.}
$12/30=2/5=40/100$. Cela représente $\res{40\,\%}$ de la classe.
\end{corrige}\end{blocLivre}
\begin{blocLivre}{cor-11-5}\begin{corrige}{5}{Un tableau à reconnaître}\demarche{Comparer les quotients correspondants.}
Les quotients du prix par le nombre de lots valent $4/1=8/2=12/3=4$. Le prix s’obtient en multipliant le nombre de lots par 4 : $\res{\text{oui}}$, le prix est proportionnel au nombre de lots, avec un coefficient de 4 euros par lot.
\end{corrige}\end{blocLivre}
\begin{blocLivre}{cor-11-6}\begin{corrige}{6}{Une recette pour six}\demarche{Employer le même coefficient pour chaque ingrédient.}
Le coefficient est $6/4=1{,}5$. Il faut $\res{375\ \mathrm g}$ de farine et $\res{30\ \mathrm{cL}}$ de lait.
\end{corrige}\end{blocLivre}
\begin{blocLivre}{cor-11-7}\begin{corrige}{7}{Une remise}\demarche{Calculer la part puis la soustraire.}
La réduction est $(80:100)\times15=12$ euros. On paie $80-12=\res{68}$ euros. 15 \% désigne ici la réduction, pas le prix final.
\end{corrige}\end{blocLivre}
\begin{blocLivre}{cor-11-8}\begin{corrige}{8}{Une échelle de plan}\demarche{Convertir avant d’appliquer l’échelle.}
$4{,}8$ m $=480$ cm. La longueur dessinée est $480:100=\res{4{,}8\ \mathrm{cm}}$.
\label{fig-cor-11-8}\figureGeo[corrige]{prop-8}
\end{corrige}\end{blocLivre}
\begin{blocLivre}{cor-11-9}\begin{corrige}{9}{Vitesse moyenne}\demarche{Exprimer la durée en heures décimales.}
$1$ h $30=1{,}5$ h. La vitesse moyenne est $27:1{,}5=\res{18\ \mathrm{km/h}}$. $1{,}30$ heure ne représente pas 1 h 30.
\end{corrige}\end{blocLivre}
\begin{blocLivre}{cor-11-10}\begin{corrige}{10}{Une situation non proportionnelle}\demarche{Tester un doublement.}
Un colis coûte 6 euros, deux coûtent 8 euros. Doubler le nombre ne double pas le prix. La situation est $\res{\text{non proportionnelle}}$.
\end{corrige}\end{blocLivre}
\begin{blocLivre}{cor-11-11}\begin{corrige}{11}{Lire des points proportionnels}\demarche{Lire les axes puis revenir à l’unité.}
Le point d’abscisse 3 a pour ordonnée 7,5 euros. Un ticket coûte $7{,}5:3=\res{2{,}50}$ euros. Les points et l’origine sont alignés.
\label{fig-cor-11-11}\figureGeo[corrige]{prop-11}
\end{corrige}\end{blocLivre}
\begin{blocLivre}{cor-11-12}\begin{corrige}{12}{Retrouver un effectif}\demarche{Retrouver une part repère.}
10 \% correspondent à $18:3=6$ personnes. 100 \% correspondent à $6\times10=\res{60}$ personnes. Contrôle : 30 \% de 60 vaut 18.
\end{corrige}\end{blocLivre}
\begin{blocLivre}{cor-11-13}\begin{corrige}{13}{Comparer deux promotions}\demarche{Comparer les prix finaux.}
A retire 12 euros et fait payer 48 euros. B fait payer 50 euros. L’offre $\res A$ est moins chère de 2 euros.
\end{corrige}\end{blocLivre}
\begin{blocLivre}{cor-11-14}\begin{corrige}{14}{Deux trajets}\demarche{Diviser les totaux.}
Distance totale : $12+18=30$ km. Durée totale : $0{,}5+1{,}5=2$ h. Vitesse moyenne : $30:2=\res{15\ \mathrm{km/h}}$. La moyenne simple de 24 et 12 ne convient pas, car les durées diffèrent.
\end{corrige}\end{blocLivre}
\begin{blocLivre}{cor-11-15}\begin{corrige}{15}{Une carte et un détour}\demarche{Mettre les unités au même niveau puis appliquer la proportion.}
$6{,}4\times25000=160000$ cm, soit 1,6 km. Le détour vaut $15\,\%$ de 1,6 km, soit 0,24 km. La distance totale est $\res{1{,}84\ \mathrm{km}}$.
\label{fig-cor-11-15}\figureGeo[corrige]{prop-15}
\end{corrige}\end{blocLivre}
